\chapter{Real-бимодуль пространства стрелок и относительный семейный quotient}
\label{ch:v7-real-arrow-bimodule-forest-quotient-gate}

\section{Два разных вопроса вложения}

Предыдущий гейт получил единый Hodge-функционал на пространстве всех
одиннадцати новых стрелок. Теперь необходимо различить два утверждения:
существование одной Real-суперсвязности на модуле стрелок и её происхождение
как обычной внутренней флуктуации конечного физического оператора Дирака.
Первое утверждение оказывается положительным, второе --- отрицательным.

Пусть
\begin{equation}
 \mathcal E_{\rm new}
 =\bigoplus_{e\in E_{\rm new}}
 \operatorname{Hom}(\mathcal H_{s(e)},\mathcal H_{t(e)}),
 \qquad |E_{\rm new}|=11.
 \label{eq:v7-real-arrow-edge-module}
\end{equation}
Проекторы $P_+$ и $P_-$ являются суммами тождественных операторов на
полных калибровочно-ковариантных слагаемых. Поэтому они коммутируют с
индуцированным действием конечной алгебры на модуле стрелок.

\section{Одна Real-суперсвязность пространства стрелок}

На двухступенном модуле
\begin{equation}
 \mathcal K_E=\mathcal E_{\rm new}^{0}\oplus
 \mathcal E_{\rm new}^{1},
 \qquad
 \chi_E=\operatorname{diag}(-I,I)
 \label{eq:v7-real-arrow-two-term-module}
\end{equation}
действие алгебры одинаково на обеих ступенях:
\begin{equation}
 \rho_E(a)=\operatorname{diag}(\rho_{\rm arr}(a),
                                \rho_{\rm arr}(a)).
 \label{eq:v7-real-arrow-algebra-action}
\end{equation}
Фоновый оператор предыдущего гейта удовлетворяет
\begin{equation}
 \delta_E=\begin{pmatrix}0&P_+\\P_-&0\end{pmatrix},
 \qquad
 \delta_E^2=0,
 \qquad
 [\delta_E,\rho_E(a)]=[d_Z,\rho_E(a)]=0.
 \label{eq:v7-real-arrow-background-bimodule}
\end{equation}
Следовательно, он является нечётным бимодульным эндоморфизмом. Вместе с
динамическим $d_Z$ он определяет один колчанный суперсвязностный объект
\begin{equation}
 \mathbb A_E=\nabla_E+d_Z+\mu\delta_E,
 \qquad
 \mathfrak m_E=[d_Z,d_Z^\dagger]
 +\mu^2[\delta_E,\delta_E^\dagger].
 \label{eq:v7-real-arrow-superconnection}
\end{equation}
Здесь два цвета являются двумя стрелками одного представления, поэтому их
вклады складываются в одном отображении момента. Это стандартная структура
пространства представлений колчана и нормы отображения момента
\cite{v7-real-arrow-harada}. Нечётная эндоморфная часть согласуется с
понятием суперсвязности Quillen \cite{v7-real-arrow-quillen}.

Real-оператор меняет ориентацию стрелки и обменивает ступени. Для формального
ориентационного удвоения выполняется
\begin{equation}
 J_E\chi_EJ_E^{-1}=-\chi_E,
 \qquad
 J_E\delta_EJ_E^{-1}=\delta_E^\dagger.
 \label{eq:v7-real-arrow-real-structure}
\end{equation}

\section{Почему это не обычная внутренняя флуктуация}

Бимодульная линейность обоих нечётных операторов, достаточная для
вспомогательной суперсвязности,
одновременно создаёт препятствие для стандартного спектрального чтения.
Если использовать этот оператор как конечный $D_E$, то
\begin{equation}
 \Omega^1_{D_E}(\mathcal A_F)
 =\operatorname{span}\{a[D_E,b]:a,b\in\mathcal A_F\}=0.
 \label{eq:v7-real-arrow-inner-oneform-zero}
\end{equation}
Поэтому динамическое поле $Z$ не получается его обычной внутренней
флуктуацией. Носитель остаётся Real-эквивариантным соответствием пространства
стрелок, а не новым фермионным конечным спектральным тройным объектом.
Классификация физических конечных троек требует оператор на фермионном
бимодуле и нетривиальные допустимые одноформы \cite{v7-real-arrow-krajewski}.

\section{Лес новых рёбер}

Вакуумная опора нового Hodge-функционала равна
\begin{equation}
 E_*=\{L_LY_R,Q_LY_R,X_LX_R,X_Le_R,X_Lu_R,Y_LY_R\}.
 \label{eq:v7-real-arrow-selected-edges}
\end{equation}
На девяти вершинах, включая изолированную относительно $E_*$ вершину $d_R$,
матрица инцидентности $B_*$ имеет
\begin{equation}
 |V|=9,\qquad |E_*|=6,\qquad c_*=3,\qquad
 \rank B_*=6,\qquad |E_*|-|V|+c_*=0.
 \label{eq:v7-real-arrow-selected-forest}
\end{equation}
То есть одни новые ненулевые рёбра образуют лес. Если бы старый фон можно
было вращать независимо вместе с ними, их шесть унитарных ориентаций не
содержали бы циклической голономии.

\section{Возврат замороженного фона \texorpdfstring{$H_{15}$}{H15}}

Однако родитель является относительным к трём уже существующим рёбрам
\begin{equation}
 E_0=\{L_Le_R,Q_Ld_R,Q_Lu_R\}.
 \label{eq:v7-real-arrow-baseline-edges}
\end{equation}
В текущем семейно-слепом подъёме их унитарные полярные ориентации равны
$I_3$. Все последующие размерности quotient относятся именно к этому
замороженному представителю; после вывода неравных семейных сингулярных
чисел стабилизатор необходимо вычислить повторно.
Полный ненулевой граф имеет девять рёбер, связан и содержит ровно один цикл:
\begin{equation}
 |E_0\cup E_*|=9,\qquad c_{\rm full}=1,\qquad
 \rank B_{\rm full}=8,\qquad
 b_1=9-9+1=1.
 \label{eq:v7-real-arrow-full-cycle-rank}
\end{equation}
Именно этот цикл совпадает с ранее найденным $H_{15}$-укоренённым маршрутом;
две вектороподобные массы являются древесными ветвями.

Пусть $U_v\in U(3)$ меняет семейный кадр в вершине. На единичном вакууме
линеаризованное действие равно
\begin{equation}
 \delta Z_e=\xi_{t(e)}-\xi_{s(e)},
 \qquad \xi_v\in\mathfrak u(3).
 \label{eq:v7-real-arrow-frame-action}
\end{equation}
Сохранение замороженных старых рёбер требует
\begin{equation}
 B_0^T\xi=0,\qquad
 \dim_{\mathbb R}\ker(B_0^T\otimes I_{\mathfrak u(3)})
 =(9-3)\cdot9=54.
 \label{eq:v7-real-arrow-baseline-frame-kernel}
\end{equation}
Ограничение действия этих кадров на новые рёбра имеет ранг
\begin{equation}
 \rank_{\mathbb R}
 \left.(B_*^T\otimes I_{\mathfrak u(3)})\right|_{\ker B_0^T}
 =5\cdot9=45.
 \label{eq:v7-real-arrow-relative-frame-rank}
\end{equation}
Поэтому из $54$ нулевых направлений семейного Hodge-гессиана снимаются
только $45$:
\begin{equation}
 \dim_{\mathbb R}\mathcal M_{\rm orient}^{\rm rel}
 =54-45=9=\dim_{\mathbb R}U(3).
 \label{eq:v7-real-arrow-relative-moduli}
\end{equation}

Оставшийся объект можно представить голономией единственного цикла
\begin{equation}
 W_C=Z_{e_1}^{\epsilon_1}Z_{e_2}^{\epsilon_2}\cdots
 Z_{e_6}^{\epsilon_6}\in U(3),
 \qquad \epsilon_j\in\{1,-1\},
 \label{eq:v7-real-arrow-cycle-holonomy}
\end{equation}
которая при вершинной смене кадров преобразуется сопряжением
\begin{equation}
 W_C\longmapsto U_{v_0}W_CU_{v_0}^{-1}.
 \label{eq:v7-real-arrow-holonomy-conjugation}
\end{equation}
Текущий функционал не зависит от $W_C$ и потому не выбирает его спектр.
Девятимерность в \eqref{eq:v7-real-arrow-relative-moduli} является
линеаризованной размерностью в единичной точке, где сопряжение имеет
неподвижную точку. Нелинейный кадровый quotient равен
\begin{equation}
 \mathcal M_C=U(3)/\operatorname{Ad}U(3),
 \qquad \dim_{\mathbb R}\mathcal M_C^{\rm generic}=3,
 \label{eq:v7-real-arrow-conjugacy-quotient}
\end{equation}
и описывается тремя собственными фазами. Даже эти три классовых инварианта
ещё не являются выводом CKM или PMNS: их связь с наблюдаемыми не установлена.

\begin{theorem}[Real-подъём и частичный семейный quotient]
Два цвета Hodge-родителя образуют один нечётный Real-эквивариантный
суперсвязностный объект на удвоенном модуле полного пространства стрелок.
Его проекторы бимодульно линейны, нильпотентны в требуемой композиции и не
добавляют физические фермионные вершины. После учёта замороженного $H_{15}$
вершинные семейные кадры снимают $45$ из $54$ нулевых мод. Оставшиеся девять
линейных направлений являются касательным конусом голономии единственного
$H_{15}$-укоренённого цикла; её нелинейный quotient имеет три собственные
фазы.
\end{theorem}

\begin{nogocriterion}[Два запрета преждевременного замыкания]
Нельзя называть модуль стрелок стандартной конечной спектральной тройкой:
его бимодульный нечётный оператор имеет нулевой модуль обычных внутренних
одноформ. Нельзя также удалить все $54$ семейные нулевые моды полным
$U(3)^9$-вращением, поскольку это одновременно вращает замороженный фон
$H_{15}$; относительный quotient оставляет один $U(3)$-цикл.
Нельзя переносить число $45$ на ещё не выведенные неравные семейные матрицы
$H_{15}$ без нового вычисления их стабилизатора.
\end{nogocriterion}

\begin{protocol}
\begin{enumerate}
 \item один Real-модуль двух цветов стрелок: \textbf{получен};
 \item бимодульный коммутатор $[\delta_E,\rho_E(a)]$: \textbf{$0$};
 \item вспомогательное условие первого порядка: \textbf{тривиально выполнено};
 \item обычные внутренние одноформы этого оператора: \textbf{$0$};
 \item стандартная физическая внутренняя флуктуация: \textbf{не получена};
 \item новых физических фермионов: \textbf{$0$};
 \item цикл-ранг новых рёбер: \textbf{$0$};
 \item цикл-ранг вместе с $H_{15}$: \textbf{$1$};
 \item семейных нулевых мод до quotient: \textbf{$54$};
 \item вертикальных кадровых мод: \textbf{$45$};
 \item относительных линейных циклических мод в единице: \textbf{$9$};
 \item нелинейных классовых параметров общего $U(3)$-цикла: \textbf{$3$};
 \item вывод CKM/PMNS: \textbf{нет};
 \item масштаб $\mu$: \textbf{не выведен};
 \item итог: \textbf{положительное Real-соответствие, физическая внутренняя
 флуктуация закрыта, один циклический модуль открыт}.
\end{enumerate}
\end{protocol}

Машинный сертификат сохранён в
\path{s2t_v7_real_arrow_bimodule_forest_quotient_gate_results.json}.

\begin{thebibliography}{9}
\bibitem{v7-real-arrow-harada}
M.~Harada, G.~Wilkin,
\emph{Morse theory of the moment map for representations of quivers},
Geometriae Dedicata 150 (2011), 307--353; arXiv:0807.4734.
\bibitem{v7-real-arrow-quillen}
D.~Quillen,
\emph{Superconnections and the Chern character},
Topology 24 (1985), 89--95.
\bibitem{v7-real-arrow-krajewski}
T.~Krajewski,
\emph{Classification of Finite Spectral Triples},
Journal of Geometry and Physics 28 (1998), 1--30;
arXiv:hep-th/9701081.
\end{thebibliography}